\chapter{Dimensionality Reduction}
\label{ch:dimensionality-reduction}

\epigraph{``The curse of dimensionality is not about computation. It is about geometry. In high dimensions, all points are far apart.''}{}

\learninggoal{
	\item Explain the curse of dimensionality and why high-dimensional data is problematic.
	\item Derive principal component analysis from the variance-maximisation perspective.
	\item Understand the SVD formulation of PCA and its connection to feature engineering.
	\item Compare linear (PCA) and non-linear (t-SNE, UMAP) dimensionality reduction.
	\item Apply PCA for noise reduction, visualization, and pre-processing.
}

\section{The Curse of Dimensionality}

As dimensionality $d$ increases, the volume of space grows exponentially. Data becomes sparse, distances converge, and models require exponentially more samples.

\begin{example}[Concentration of Distances]
	For $n$ points uniformly distributed in $[0, 1]^d$, the ratio of the nearest-neighbor distance to the farthest-neighbor distance approaches 1 as $d$ grows:
	\[
		\frac{\text{dist}_{\text{nearest}}}{\text{dist}_{\text{farthest}}} \to 1
	\]
	In high dimensions, the concept of ``nearest neighbor'' loses meaning.
\end{example}

Consequences:
\begin{itemize}
	\item KNN classifiers perform poorly above $d \approx 20$ without feature selection.
	\item Density estimation requires $n \propto 2^d$ samples.
	\item Overfitting risk increases as $d$ approaches $n$.
\end{itemize}

\section{Principal Component Analysis}

PCA is the most widely used linear dimensionality reduction method. It finds a sequence of orthogonal directions that maximize the variance of the projected data.

\begin{definition}[PCA Objective]
	Find $\mathbf{w}_1$ (first principal component) that maximizes the variance of $\mathbf{X}\mathbf{w}_1$:
	\[
		\mathbf{w}_1 = \arg\max_{\|\mathbf{w}\|=1} \Var(\mathbf{X}\mathbf{w}) = \arg\max_{\|\mathbf{w}\|=1} \mathbf{w}^\top \boldsymbol{\Sigma} \mathbf{w}
	\]
	where $\boldsymbol{\Sigma} = \frac{1}{n} \mathbf{X}^\top \mathbf{X}$ is the covariance matrix (assuming centred data).
\end{definition}

\subsection{Eigenvalue Solution}

The solution is the eigenvector of $\boldsymbol{\Sigma}$ with the largest eigenvalue:

\[
	\boldsymbol{\Sigma} \mathbf{w}_1 = \lambda_1 \mathbf{w}_1
\]

Subsequent components $\mathbf{w}_2, \dots, \mathbf{w}_k$ are the next eigenvectors, orthogonal to all previous ones.

\subsection{PCA via SVD}

PCA is more numerically stable when computed via the SVD of the data matrix:

\[
	\mathbf{X} = \mathbf{U} \boldsymbol{\Sigma} \mathbf{V}^\top
\]

\begin{itemize}
	\item Columns of $\mathbf{V}$: principal component directions (same as eigenvectors of $\mathbf{X}^\top \mathbf{X}$).
	\item Singular values $\sigma_j$: related to eigenvalues by $\lambda_j = \sigma_j^2 / n$.
	\item Principal component scores: $\mathbf{Z} = \mathbf{X} \mathbf{V}$.
	\item Variance explained by component $j$: $\sigma_j^2 / \sum_{i=1}^d \sigma_i^2$.
\end{itemize}

\subsection{Choosing the Number of Components}

A scree plot shows the variance per component. Common criteria:

\begin{itemize}
	\item \textbf{Kaiser rule:} keep components with eigenvalue $> 1$ (for correlation matrix).
	\item \textbf{Cumulative variance:} keep enough components to explain $90\%$ or $95\%$ of variance.
	\item \textbf{Elbow method:} look for the knee in the scree plot.
\end{itemize}

\subsection{PCA as Feature Engineering}

PCA creates new features (principal components) that are:
\begin{itemize}
	\item Uncorrelated (orthogonal).
	\item Ordered by variance (information content).
	\item Linear combinations of original features.
	\item Useful for visualization (first 2--3 components).
	\item Effective denoising (discard low-variance components).
\end{itemize}

\begin{principle}[PCA for Linear Models]
	Since PCA components are uncorrelated, linear models fit on PCA-transformed data have no multicollinearity. This stabilizes coefficient estimates and can improve generalization when $d \gg n$.
\end{principle}

\section{Factor Analysis}

Factor analysis is similar to PCA but with a different goal: model the covariance structure.

\[
	\mathbf{x} = \boldsymbol{\mu} + \mathbf{L} \mathbf{f} + \boldsymbol{\epsilon}
\]

\begin{itemize}
	\item $\mathbf{f} \sim \N(0, \mathbf{I})$: latent factors ($k \ll d$).
	\item $\mathbf{L} \in \R^{d \times k}$: factor loadings (how each feature relates to each factor).
	\item $\boldsymbol{\epsilon} \sim \N(0, \boldsymbol{\Psi})$: unique variances (diagonal).
\end{itemize}

Unlike PCA, factor analysis explicitly models measurement error and assumes the observed correlations arise from a small number of common causes.

\section{t-Distributed Stochastic Neighbor Embedding}

t-SNE is a non-linear method primarily used for visualization in 2D or 3D.

\begin{itemize}
	\item \textbf{Local structure preservation:} points close in high-dimensional space remain close in low-dimensional space.
	\item \textbf{Probability-based:} convert distances to conditional probabilities $p_{j|i}$, then minimize KL divergence between high-dim and low-dim distributions.
	\item \textbf{Heavy-tailed distribution:} uses a Student-$t$ distribution in the low-dimensional space to alleviate the crowding problem.
\end{itemize}

\begin{lstlisting}[language=Python, caption=t-SNE for visualization]
from sklearn.manifold import TSNE

tsne = TSNE(n_components=2, perplexity=30, random_state=0)
X_2d = tsne.fit_transform(X)
\end{lstlisting}

Key hyperparameter: \texttt{perplexity} (roughly the effective number of neighbors). Typical range: 5--50.

\subsection{Limitations of t-SNE}

\begin{itemize}
	\item Non-deterministic: different runs give different results.
	\item Global structure (distances between clusters) is not preserved.
	\item Computationally expensive: $O(n^2)$.
	\item Cannot embed new points (transductive, not inductive).
\end{itemize}

\section{UMAP}

Uniform Manifold Approximation and Projection (UMAP) addresses some of t-SNE's limitations:

\begin{itemize}
	\item Faster: $O(n \log n)$ via approximate nearest neighbors.
	\item Better global structure preservation.
	\item Supports embedding new points after training.
	\item Based on manifold theory and fuzzy simplicial sets.
\end{itemize}

UMAP is preferred over t-SNE for large datasets and when reproducibility matters.

\section{Autoencoders}

Neural network-based dimensionality reduction. An autoencoder learns to reconstruct its input through a bottleneck layer of dimension $k \ll d$:

\[
	\mathbf{h} = f(\mathbf{W}_e \mathbf{x} + \mathbf{b}_e)
\]
\[
	\hat{\mathbf{x}} = g(\mathbf{W}_d \mathbf{h} + \mathbf{b}_d)
\]

The bottleneck $\mathbf{h}$ is a non-linear low-dimensional representation. Autoencoders can capture manifolds that linear methods (PCA) cannot.

\section{Choosing a Method}

\begin{longtable}{p{2.5cm}p{4cm}p{4.5cm}p{3cm}}
	\toprule
	\textbf{Method} & \textbf{Type}          & \textbf{Best for}         & \textbf{Complexity}   \\
	\midrule
	PCA             & Linear                 & Pre-processing, denoising & $O(\min(nd^2, dn^2))$ \\
	Factor analysis & Linear (probabilistic) & Covariance modelling      & $O(ndk)$              \\
	t-SNE           & Non-linear             & Visualization (2D/3D)     & $O(n^2)$              \\
	UMAP            & Non-linear             & Visualization, large data & $O(n \log n)$         \\
	Autoencoder     & Non-linear (neural)    & Feature learning          & $O(n d h)$            \\
	\bottomrule
\end{longtable}

\section{Exercises}

\begin{exercise}
	Generate 1000 points uniformly in $[0, 1]^d$ for $d = 2, 10, 50, 100$. For each dimension, compute the ratio of maximum to minimum pairwise distance. At what $d$ does the ratio approach 1?
\end{exercise}

\begin{exercise}
	Apply PCA to the Iris dataset. Compute the variance explained by each component. How many components capture $95\%$ of the variance? Visualize the data projected onto the first two principal components, colored by species.
\end{exercise}

\begin{exercise}
	Compare PCA and t-SNE on the MNIST dataset. Plot the 2D embeddings side by side. Which method produces better separation of digit classes? Which method preserves the global structure (e.g., the fact that 3 and 8 are more similar than 3 and 0)?
\end{exercise}

\begin{exercise}
	Use PCA as a pre-processing step for linear regression on a dataset with $d > n$. Compare test error with and without PCA. How many components minimize test error? Why does PCA help here?
\end{exercise}
